\subsection{The Value of Adaptivity at a Given Population}\label{app:population}
Fix a menu $\A$ and a law $P$ on $\X\times\{-1,1\}$. A \emph{design} is a
probability vector $(q_\iota)$ on $\A$ with tables $g_\iota(\ell,y)$ that
satisfy $\sum_\iota g_\iota(\ell_\iota(x),y)=t(x,y)=yd(x)$ for all $(x,y)$.
Its score is $H=g_\iota(\ell_\iota(X),Y)/q_\iota$ with $\iota\sim q$ drawn
independently of the record, which is unbiased for $t$ at every profile and
label; conversely every function $H$ of $(\iota,\ell_\iota(X),Y)$ with
$\E[H\mid X,Y]=t$ arises in this way. $W_\A(P)$ is the infimum of
$\Var_{P,q}(H)$ over designs, $M_\A(P)=W_\A(P)+\Delta(P)^2$ the infimum of
the second moment, and $\sigma_P^2=\Var_P(t)$. At a full-support law,
$W_\A(P)$ is also the local asymptotic minimax risk over arbitrary estimators
(Proposition~\ref{prop:local}), so ratios of $W$ are ratios of the numbers of
records that two menus need. For a deterministic depth-$b$ tree $\tau$ that
computes $d$, $S_\tau(x)$ is its padded path set and $\mathcal S_\tau$ the
collection of these sets (Appendix~\ref{app:universal}).

\begin{theorem}[Records needed at a population, full statement of Theorems~\ref{thm:universal}(ii), \ref{thm:rounds}(a) and~\ref{thm:bracket}]\label{thm:bracketfull}
Let $d$ be a nonzero contrast of degree at most $b$ on $\{-1,1\}^m$ and $P$ any
law on $\X\times\{-1,1\}$.

\textup{(a) Uniform gap.} For $1\le r\le b$,
\[
 W_{\rm branch}(P)\ \le\ W_r(P)\ \le\ \min\Big\{2^{b-r},\binom{m-r}{b-r}\Big\}\,W_{\rm branch}(P),
\]
and in particular $W_{\rm fixed}(P)\le K_{m,b}W_{\rm branch}(P)$.

\textup{(b) Witness bound.} For every menu $\A$, every constant $c$ and
every function $h$ on $\X\times\{-1,1\}$ with $\E_P[(t-c)h]\ne0$,
\[
 W_\A(P)+(\Delta(P)-c)^2\ \ge\ \frac{\E_P[(t-c)h]^2}{\max_{\iota\in\A}\E_P\big[\E_P[h\mid\ell_\iota(X),Y]^2\big]}.
\]
With $c=0$ this bounds the second moment $M_\A(P)$, and with $c=\Delta(P)$ it
bounds the variance $W_\A(P)$.

\textup{(c) Bracket.} If a depth-$b$ tree $\tau$ computes $d$ and
$\sigma_P^2>0$, then $W_{\rm branch}(P)=\sigma_P^2$ and
\[
 \frac1{\rho_P^2}\ \le\ \frac{W_{\rm fixed}(P)}{W_{\rm branch}(P)}\ \le\ N_P(\tau)\ \le\ |\mathcal S_\tau|\ \le\ K_{m,b},
\]
where
\[
 \rho_P^2=\max_{|T|=b}\frac{\Var_P\big(\E_P[t\mid X_T,Y]\big)}{\Var_P(t)},
 \qquad
 N_P(\tau)=\Big(\sum_{S\in\mathcal S_\tau}v_S^{1/2}\Big)^2,
 \quad v_S=\frac{\E_P\big[(t-\Delta(P))^2\mathbf 1\{S_\tau(X)=S\}\big]}{\sigma_P^2}.
\]
More generally $W_{\rm fixed}(P)\le\big(\sum_S\E_P[(t-c)^2\mathbf 1\{S_\tau(X)=S\}]^{1/2}\big)^2-(\Delta(P)-c)^2$
for every constant $c$.

\textup{(d)} For the address contrast with $\kappa=b-1$ under the uniform law
with independent uniform labels, $1/\rho_P^2=N_P(\tau)=K_{m,b}=D$, so all
inequalities in \textup{(c)} are equalities.
\end{theorem}

The quantity $\rho_P^2$ is the largest fraction of the variance of the target
that a single form explains, a population $R^2$. The weights $v_S$ sum to
one, and $N_P(\tau)=\exp H_{1/2}(v)$ is the exponential of their R\'enyi
entropy of order $1/2$: an effective number of question sets that the tree
uses at $P$. Adaptivity can be valuable only when the tree's questions diverge
over many well-populated paths, and it is valuable when in addition no fixed
set of $b$ questions explains most of the variance of the target; correlated
features and shared coordinates raise $\rho_P^2$.

\begin{proof}
\emph{(a)} The left inequality is menu inclusion. For the right one fix a
branching design: trees $\tau$ with probabilities $q_\tau>0$ and tables
$g_\tau(\ell,y)$ with $\sum_\tau g_\tau(\ell_\tau(x),y)=t(x,y)$, and write
$H_\tau(x,y)=g_\tau(\ell_\tau(x),y)/q_\tau$, so that
$\sum_\tau q_\tau H_\tau=t$ pointwise. Pad every tree as in
Appendix~\ref{app:universal}; padding refines the transcript, so $H_\tau$
is still a function of the padded leaf and the label. For a node $v$ of
$\tau$ at depth $r-1$ (the root if $r=1$), the padded subtree below $v$ has
depth $b-r+1$; let $\mathcal S_{\tau,v}$ be the collection of coordinate sets
of its root-to-leaf paths and $L_{\tau,v}=|\mathcal S_{\tau,v}|$. These sets
contain the coordinate asked at $v$ and are determined by the first $b-r$
answers below $v$, so $L_{\tau,v}\le\min\{2^{b-r},\binom{m-r}{b-r}\}=:K_r$.
Fix a constant $c$ and consider the following randomized $r$-round interview
with at most $b$ queries. Draw $\tau\sim q$ and, independently of the record,
one set $S_v$ uniformly from $\mathcal S_{\tau,v}$ for every node $v$ of
depth $r-1$. In rounds $1,\ldots,r-1$ ask one question per round and follow
$\tau$ to a node $v$; in round $r$ ask the coordinates of $S_v$. The path of
$x$ below $v$ can be followed to a leaf exactly when its coordinate set equals
$S_v$, which has conditional probability $1/L_{\tau,v}$ given $(x,y,\tau)$,
and this event is measurable from the transcript. Return
\[
 \widetilde H=c+L_{\tau,v}\big(H_\tau(x,y)-c\big)\ \text{on completion},\qquad \widetilde H=c\ \text{otherwise}.
\]
The pair $(\tau,(S_v)_v)$ indexes a deterministic $r$-round interview, so
$\widetilde H$ is a score of a design on $\A_r(b)$. Conditionally on $(x,y)$
and $\tau$, its mean is $H_\tau(x,y)$ and
$\E[(\widetilde H-c)^2\mid x,y,\tau]=L_{\tau,v}(H_\tau(x,y)-c)^2\le K_r(H_\tau(x,y)-c)^2$.
Hence $\E[\widetilde H\mid X,Y]=t$, and with $H=H_\tau$ for $\tau\sim q$,
\[
 \Var(\widetilde H)=\E[(\widetilde H-c)^2]-(\Delta(P)-c)^2
 \le K_r\,\E[(H-c)^2]-(\Delta(P)-c)^2 .
\]
Choose $c=\Delta(P)$: then $\E[(H-c)^2]=\Var(H)$ and
$\Var(\widetilde H)\le K_r\Var(H)$. The constant $c$ is part of the score, not
an estimate: the design is pointwise unbiased for every value of $c$, and
$W_r(P)$ is by definition an infimum over designs that may be tuned to $P$.
Taking the infimum over branching designs proves (a).

\emph{(b)} Let $(q,g)$ be a design on $\A$ with score $H$. The tables
$g_\iota-cq_\iota$ sum to $t-c$, and
$\E[(H-c)^2]=\sum_\iota\|g_\iota-cq_\iota\|_P^2/q_\iota\ge(\sum_\iota\|g_\iota-cq_\iota\|_P)^2$
by Cauchy--Schwarz in $q$. For every $h$,
\[
 \E_P[(t-c)h]=\sum_\iota\E_P\big[(g_\iota-cq_\iota)\,\E_P[h\mid\ell_\iota(X),Y]\big]
 \le\Big(\sum_\iota\|g_\iota-cq_\iota\|_P\Big)\max_\iota\big\|\E_P[h\mid\ell_\iota(X),Y]\big\|_P,
\]
because $g_\iota-cq_\iota$ is a function of $(\ell_\iota(X),Y)$. Square, use
$\E[(H-c)^2]=\Var(H)+(\Delta(P)-c)^2$ and take the infimum over designs.

\emph{(c)} Following $\tau$ returns $t$ itself, and every unbiased score has
$\Var(H)\ge\Var(\E[H\mid X,Y])=\sigma_P^2$, so $W_{\rm branch}(P)=\sigma_P^2$.
For the lower bound apply (b) to the fixed forms with $c=\Delta(P)$ and
$h=t-\Delta(P)$: the numerator is $\sigma_P^4$ and
$\E_P[\E_P[t-\Delta\mid X_T,Y]^2]=\Var_P(\E_P[t\mid X_T,Y])$ because
$\E_P[t\mid X_T,Y]$ has mean $\Delta(P)$. For the upper bound choose $q_S>0$
on $\mathcal S_\tau$, draw $S$ before the person answers and return
$c+(t(X,Y)-c)/q_S$ if the tree's path completes inside $S$ and $c$ otherwise.
Completion occurs exactly when $S=S_\tau(X)$
(Appendix~\ref{app:universal}), so the score is unbiased at every profile and
label, and $\E[(H-c)^2]=\sum_Sq_S^{-1}\E_P[(t-c)^2\mathbf 1\{S_\tau(X)=S\}]$.
Minimizing over the simplex gives
$q_S\propto\E_P[(t-c)^2\mathbf 1\{S_\tau=S\}]^{1/2}$ and the value
$(\sum_S\E_P[(t-c)^2\mathbf 1\{S_\tau=S\}]^{1/2})^2$; sets of zero weight are
handled by letting their probability tend to zero. Subtracting
$(\Delta(P)-c)^2$ gives the general statement, and $c=\Delta(P)$ gives
$N_P(\tau)\sigma_P^2$. Finally $(\sum_Sv_S^{1/2})^2\le|\mathcal S_\tau|\sum_Sv_S=|\mathcal S_\tau|$
by Cauchy--Schwarz, and $|\mathcal S_\tau|\le K_{m,b}$.

\emph{(d)} Here $\Delta=0$, $t^2=1$ and the $D$ path sets are equally likely,
so $N_P=D=K_{m,b}$. With independent labels,
$\E[t\mid X_T,Y]=Y\E[d\mid X_T]$, whose second moment is at most $2^{-\kappa}=1/D$
by~\eqref{eq:addressinformation}, so $1/\rho_P^2\ge D$.
\end{proof}

\paragraph{Relation to the uncentered bracket.}
Setting $c=0$ in (b) and (c) gives the second-moment bracket
$\E_P[t^2]/\Gamma_P\le M_{\rm fixed}(P)\le N^0_P(\tau)\E_P[t^2]$ with
$\Gamma_P=\max_T\E_P[\E_P[t\mid X_T,Y]^2]/\E_P[t^2]$ and $N_P^0$ computed
from the weights $\E_P[t^2\mathbf 1\{S_\tau=S\}]$. Subtracting $\Delta^2$
turns it into a bracket for $W_{\rm fixed}/W_{\rm branch}$ with lower end
$(\E_P[t^2]/\Gamma_P-\Delta^2)/\sigma_P^2$ and upper end
$(N_P^0\E_P[t^2]-\Delta^2)/\sigma_P^2$. The centered lower end is never
smaller: writing $s=\E_P[t^2]$ and $e=\Gamma_Ps$, the claim
$(s-\Delta^2)^2\ge(s^2/e-\Delta^2)(e-\Delta^2)$ reduces to $s^2/e+e\ge2s$. The
uncentered upper end is unbounded as $\Delta^2\to\E_P[t^2]$, that is for
accurate rules, whereas the centered one never exceeds $K_{m,b}$; for the
worked example of Section~\ref{sec:setup} with uniform profiles and a label
that agrees with $d$ with probability $0.9$, the centered upper end is $2$ and
the uncentered one is $(2-0.8^2)/0.36\approx3.78$. Centering is what makes part (a) hold at every
population: without it, path guessing inflates the second moment by $K_r$
and the variance by more.

\paragraph{Label-dependent designs.}
Labels are observed before the questions are asked, so a collector may let
the interview probabilities depend on the label, as in outcome-dependent
two-phase sampling. Write $\A^{\rm lab}$ for this larger class of designs on a
menu $\A$: probability vectors $q_{\cdot\mid y}$ for $y=\pm1$ with tables that
sum to $t(\cdot,y)$ for each label.

\begin{proposition}[Label-dependent designs]\label{prop:labeldesign}
\textup{(i)} $V(d,\A)$ is unchanged if interview probabilities may depend on
the current label (Theorem~\ref{thm:asym}).
\textup{(ii)} For every law $P$, $M_{\A}^{\rm lab}(P)=\sum_{y=\pm1}M_\A(P^{(y)})$,
where $P^{(y)}=P(\cdot,y)$ is the sub-probability law of the profiles with
label $y$ and $M_\A(P^{(y)})$ is the optimum of the program for the target
$yd$ under that mass; hence $W^{\rm lab}_\A(P)=\sum_yM_\A(P^{(y)})-\Delta(P)^2\le W_\A(P)$.
\textup{(iii)} Part \textup{(a)} of Theorem~\ref{thm:bracketfull} holds for
label-dependent designs: $W^{\rm lab}_r(P)\le K_rW^{\rm lab}_{\rm branch}(P)$,
and $W^{\rm lab}_{\rm branch}(P)=\sigma_P^2$ when a depth-$b$ tree computes
$d$.
\end{proposition}
\begin{proof}
(i) The converse of Theorem~\ref{thm:asym} allows the interview to depend on
the current label, and label-independent designs attain $V$. (ii) The second
moment of a label-dependent design is
$\sum_y\sum_\iota\E_P[g_{\iota,y}(\ell_\iota(X))^2\mathbf 1\{Y=y\}]/q_{\iota\mid y}$,
and neither the constraints nor the probability vectors couple the two
labels. (iii) The simulation in the proof of
Theorem~\ref{thm:bracketfull}(a) is carried out separately for each label
with the same constant $c=\Delta(P)$; the conditional bound
$\E[(\widetilde H-c)^2\mid x,y]\le K_r\E[(H-c)^2\mid x,y]$ holds for each
$(x,y)$. The last statement is the Jensen bound.
\end{proof}
Each $M_\A(P^{(y)})$ is computed by the solver of Appendix~\ref{app:engine}
with the same certificates. Appendix~\ref{app:screendetail} reports that
label-dependent form probabilities lower the records ratio of the screening
trees by at most \nLabelMax\%, and by at most \nStratMax\% when they are also
stratified by sex.
